\section{Domänenspezifische Input-Validierung zur Krypto-Sicherung}

Ein deterministischer Blind Index entfaltet seine Schutz- und Filterfunktion nur dann fehlerfrei, wenn die zugrundeliegenden Daten in absoluter Konsistenz vorliegen. Geringfügige Abweichungen in der Formatierung (z.\,B. Leerzeichen, unterschiedliche Trennzeichen bei Datumsangaben oder uneinheitliche Groß-/Kleinschreibung) führen aufgrund der Lawineneigenschaft kryptografischer Hashfunktionen zu vollständig kollidierenden Blind Indizes. Datenqualität ist bei verschlüsselten Attributen somit direkt an die funktionale Suchbarkeit gekoppelt.

Um zu verhindern, dass malformierte oder korrupte Daten die kryptografische Pipeline im Backend blockieren oder unauffindbare Datensätze erzeugen, wurde eine strikte domänenspezifische Input-Validierung implementiert.

\subsection{Technische Umsetzung}

Die Validierungskette wurde direkt vor der Verschlüsselungs- und Hashing-Logik verankert und nutzt das \textit{Jakarta Bean Validation}-Ökosystem (Hibernate Validator):

\begin{itemize}
    \item \textbf{Postleitzahlen-Validierung (\texttt{@ValidPlz}):} Mittels einer dedizierten Validierungs-Annotation und der Definition des regulären Ausdrucks \texttt{\^{}[0-9]\{5\}\$} wird sichergestellt, dass Postleitzahlen ausnahmslos aus exakt fünf numerischen Ziffern bestehen. Syntaktische Fehler (wie \texttt{"1234"}, \texttt{"ABC12"} oder \texttt{"12345a"}) werden im Vorfeld hart abgewiesen.
    \item \textbf{Geburtsdatums-Validierung (\texttt{@ValidBirthday}):} Diese Logik erzwingt das standardisierte Datumsformat \texttt{ISO\_LOCAL\_DATE} (\texttt{YYYY-MM-DD}). Gleichzeitig wird über eine semantische Prüfung verifiziert, dass das angegebene Datum zwingend in der Vergangenheit liegt. Dies verhindert logische Artefakte und sichert eine homogene Datenbasis für die Berechnung des \texttt{birthdayHash}.
    \item \textbf{Controller-Integration:} Die Annotationen wurden deklarativ auf den Feldern der Datentransferobjekte (\texttt{PersonCreateRequest.kt}) platziert. Durch das Anfügen der \texttt{@Valid}-Annotation am Endpunkt des \texttt{PatientController.kt} prüft das Framework die Integrität der Daten, noch bevor die Request-Payload die Geschäftslogik erreicht.
\end{itemize}

\subsection{Validierung und Verifikation}

Die Effektivität der Validierungs-Pipeline wurde durch das Absetzen eines syntaktisch invaliden Datensatzes (unzureichende PLZ-Länge und Geburtsdatum in der Zukunft) überprüft:

\begin{verbatim}
curl -s -X POST http://localhost:8080/api/patients/new \
  -H "Content-Type: application/json" \
  -d '{"plz":"1234","birthday":"2090-01-01", ...}'
\end{verbatim}

Das System fängt die fehlerhaften Daten sofort ab, bricht die Krypto-Pipeline ab und antwortet mit dem HTTP-Status \texttt{400 Bad Request}. Die strukturierte JSON-Antwort maskiert interne Systemdetails und liefert dem Client präzise, domänenspezifische Fehlermeldungen:

\begin{verbatim}
"PLZ muss genau 5 Ziffern enthalten (z.B. 80331)"
"Geburtsdatum muss im Format YYYY-MM-DD vorliegen und in der Vergangenheit liegen"
\end{verbatim}

Durch diese vorgeschaltete Validierungsarchitektur wird eine kompromisslose Datenintegrität garantiert, wodurch die Zuverlässigkeit der Blind-Index-Suche im klinischen Alltag sichergestellt ist.